\documentclass{article}

\usepackage{xcolor}
\usepackage{graphicx}
\usepackage{fancyvrb}
\usepackage{listings}
\usepackage[T1]{fontenc}
\usepackage{hyperref}
\usepackage{amsmath}

\definecolor{officegreen}{rgb}{0, 0.5, 0}
\definecolor{navy}{rgb}{0, 0, 0.5}
\definecolor{linecolor}{rgb}{0.5, 0.6875, 0.6875}
\definecolor{outputcolor}{rgb}{0.375, 0.375, 0.375}

\newcommand{\id}[1]{\textcolor{black}{#1}}
\newcommand{\com}[1]{\textcolor{officegreen}{#1}}
\newcommand{\inact}[1]{\textcolor{gray}{#1}}
\newcommand{\kwd}[1]{\textcolor{navy}{#1}}
\newcommand{\num}[1]{\textcolor{officegreen}{#1}}
\newcommand{\ops}[1]{\textcolor{purple}{#1}}
\newcommand{\prep}[1]{\textcolor{purple}{#1}}
\newcommand{\str}[1]{\textcolor{olive}{#1}}
\newcommand{\lines}[1]{\textcolor{linecolor}{#1}}
\newcommand{\fsi}[1]{\textcolor{outputcolor}{#1}}
\newcommand{\omi}[1]{\textcolor{gray}{#1}}

% Overriding color and style of line numbers
\renewcommand{\theFancyVerbLine}{
\lines{\small \arabic{FancyVerbLine}:}}

\lstset{%
  backgroundcolor=\color{gray!15},
  basicstyle=\ttfamily,
  breaklines=true,
  columns=fullflexible
}

\title{{page-title}}
\date{}

\begin{document}

\maketitle




\href{https://mybinder.org/v2/gh/fsprojects/fsharp.formatting/master?filepath=codeformat.ipynb}{\begin{figure}[htbp]\centering
\includegraphics[width=1.0\textwidth]{img/badge-binder.svg}
\caption{Binder}
\end{figure}
} 
\href{https://fsprojects.github.io/FSharp.Formatting//codeformat.fsx}{\begin{figure}[htbp]\centering
\includegraphics[width=1.0\textwidth]{img/badge-script.svg}
\caption{Script}
\end{figure}
} 
\href{https://fsprojects.github.io/FSharp.Formatting//codeformat.ipynb}{\begin{figure}[htbp]\centering
\includegraphics[width=1.0\textwidth]{img/badge-notebook.svg}
\caption{Notebook}
\end{figure}
}
\section*{Code formatting}



This page demonstrates how to use \texttt{FSharp.Formatting.CodeFormat} to tokenize
F\# source code, obtain information about the source code (mainly tooltips
from the type-checker) and how to turn the code into a nicely formatted HTML.


First, we need to load the assembly and open necessary namespaces:
\begin{Verbatim}[commandchars=\\\{\}]
\kwd{open} \id{FSharp}{.}\id{Formatting}{.}\id{CodeFormat}
\kwd{open} \id{System}{.}\id{Reflection}
\end{Verbatim}

\subsection*{Starting a background agent}



The \texttt{FSharp.Formatting.CodeFormat} namespace contains \href{https://fsprojects.github.io/FSharp.Formatting/reference/fsharp-formatting-codeformat-codeformat.html}{CodeFormat} type which is the
entry point. The static method \href{https://fsprojects.github.io/FSharp.Formatting/reference/fsharp-formatting-codeformat-codeformat.html\#CreateAgent}{CodeFormat.CreateAgent} starts a background worker that
can be called to format snippets repeatedly:
\begin{Verbatim}[commandchars=\\\{\}]
\kwd{let} \id{formattingAgent} \ops{=} \ltyp{CodeFormat}{.}\id{CreateAgent}{(}{)}
\end{Verbatim}



If you want to process multiple snippets, it is a good idea to keep the
formatting agent around if possible. The agent needs to load the F\# compiler
(which needs to load various files itself) and so this takes a long time.
\subsection*{Processing F\# source}



The formatting agent provides a \href{https://fsprojects.github.io/FSharp.Formatting/reference/fsharp-formatting-codeformat-codeformatagent.html}{CodeFormatAgent.ParseAndCheckSource} method (together with an asynchronous
version for use from F\# and also a version that returns a .NET \texttt{Task} for C\#).
To call the method, we define a simple F\# code as a string:
\begin{Verbatim}[commandchars=\\\{\}]
\kwd{let} \id{source} \ops{=} \str{"""}
\str{    let hello () =}
\str{      printfn "Hello world"}
\str{  """}
\kwd{let} \id{snippets}{,} \id{errors} \ops{=} \id{formattingAgent}{.}\id{ParseAndCheckSource}{(}\str{"C:{\textbackslash}{\textbackslash}snippet.fsx"}{,} \id{source}{)}
\end{Verbatim}



When calling the method, you need to specify a file name and the actual content
of the script file. The file does not have to physically exist. It is used by the
F\# compiler to resolve relative references (e.g. \texttt{\#r}) and to automatically name
the module including all code in the file.


You can also specify additional parameters, such as \texttt{*.dll} references, by passing
a third argument with compiler options (e.g. \texttt{"-r:Foo.dll -r:Bar.dll"}).


This operation might take some time, so it is a good idea to use an asynchronous
variant of the method. It returns two arrays - the first contains F\# snippets
in the source code and the second contains any errors reported by the compiler.
A single source file can include multiple snippets using the same formatting tags
as those used on \href{http://www.fssnip.net}{fssnip.net} as documented in the
\href{http://www.fssnip.net/pages/About}{about page}.
\subsection*{Working with returned tokens}



Each returned snippet is essentially just a collection of lines and each line
consists of a sequence of tokens. The following snippet prints basic information
about the tokens of our sample snippet:
\begin{Verbatim}[commandchars=\\\{\}]
\com{// Get the first snippet and obtain list of lines}
\kwd{let} {(}{Snippet}{(}\id{title}{,} \id{lines}{)}{)} \ops{=} \id{snippets} \ops{|>} \ltyp{Seq}{.}\id{head}

\com{// Iterate over all lines and all tokens on each line}
\kwd{for} {(}{Line}{(}\id{\_}{,} \lfun{tokens}{)}{)} \kwd{in} \id{lines} \kwd{do}
  \kwd{for} \lfun{token} \kwd{in} \lfun{tokens} \kwd{do}
    \kwd{match} \lfun{token} \kwd{with}
    {|} \ltyp{TokenSpan}{.}\id{Token}{(}\lfun{kind}{,} \lfun{code}{,} \lfun{tip}{)} \kwd{->}
        \lfun{printf} \str{"}\lprf{\%s}\str{"} \lfun{code}
        \lfun{tip} \ops{|>} \ltyp{Option}{.}\id{iter} {(}\kwd{fun} \lfun{spans} \kwd{->}
          \lfun{printfn} \str{"}\lprf{\%A}\str{"} \lfun{spans}{)}
    {|} \ltyp{TokenSpan}{.}\id{Omitted} \id{\_}
    {|} \ltyp{TokenSpan}{.}\id{Output} \id{\_}
    {|} \ltyp{TokenSpan}{.}\id{Error} \id{\_} \kwd{->} {(}{)}
  \lfun{printfn} \str{""}
\end{Verbatim}



The \texttt{TokenSpan.Token} is the most important kind of token. It consists of a kind
(identifier, keyword, etc.), the original F\# code and tool tip information.
The tool tip is further formatted using a simple document format, but we simply
print the value using the F\# pretty printing, so the result looks as follows:
\begin{Verbatim}[commandchars=\\\{\}]
\kwd{let} \id{hello}{[}\id{Literal} \str{"val hello : unit -> unit"}{;} \ops{..}{.}{]} {(}{)} \ops{=}
  \id{printfn}{[}\id{Literal} \str{"val printfn : TextWriterFormat<'T> -> 'T"}{;} \ops{..}{.}{]} \str{"Hello world"}

\end{Verbatim}



The \texttt{Omitted} token is generated if you use the special \texttt{(*[omit:...]*)} command.
The \texttt{Output} token is generated if you use the \texttt{// [fsi:...]} command to format
output returned by F\# interactive. The \texttt{Error} command wraps code that should be
underlined with a red squiggle if the code contains an error.
\subsection*{Generating HTML output}



Finally, the \texttt{CodeFormat} type also includes a method \href{https://fsprojects.github.io/FSharp.Formatting/reference/fsharp-formatting-codeformat-codeformat.html}{CodeFormat.FormatHtml} that can be used
to generate nice HTML output from an F\# snippet. This is used, for example, on
\href{http://www.fssnip.net}{F\# Snippets}. The following example shows how to call it:
\begin{Verbatim}[commandchars=\\\{\}]
\kwd{let} \id{prefix} \ops{=} \str{"fst"}
\kwd{let} \id{html} \ops{=} \ltyp{CodeFormat}{.}\id{FormatHtml}{(}\id{snippets}{,} \id{prefix}{)}

\com{// Print all snippets, in case there is more of them}
\kwd{for} \lfun{snip} \kwd{in} \id{html}{.}\id{Snippets} \kwd{do}
  \lfun{printfn} \str{"}\lprf{\%s}\str{"} \lfun{snip}{.}\id{Content}

\com{// Print HTML code that is generated for ToolTips}
\lfun{printfn} \str{"}\lprf{\%s}\str{"} \id{html}{.}\id{ToolTip}
\end{Verbatim}



If the input contains multiple snippets separated using the \texttt{//[snippet:...]} comment, e.g.:
<table class="pre"><tr><td class="lines"><pre class="fssnip">
<span class="l">1: </span>
<span class="l">2: </span>
<span class="l">3: </span>
<span class="l">4: </span>
<span class="l">5: </span>
<span class="l">6: </span>
<span class="l">7: </span>
</pre>
</td>
<td class="snippet"><pre class="fssnip"><span class="c">// [snippet: First sample]</span>
<span class="i">printf</span> <span class="s">"The answer is: %A"</span> <span class="n">42</span>
<span class="c">// [/snippet]</span>
<span class="c">// [snippet: Second sample]</span>
<span class="i">printf</span> <span class="s">"Hello world!"</span>
<span class="c">// [/snippet]</span>
</pre>
</td>
</tr>
</table>


then the formatter returns multiple HTML blocks. However, the generated tool tips
are shared by all snippets (to save space) and so they are returned separately.


\end{document}